\documentclass[11pt]{article}
\usepackage[T1]{fontenc}
\usepackage[includeheadfoot,margin=0.8in,top=0.4in,bottom=0.2in]{geometry}
\usepackage[hidelinks]{hyperref}
\usepackage{lastpage}
\usepackage{mathpazo}
\usepackage{setspace}
\usepackage{fancyhdr}
\pagestyle{fancy}
\fancyhf{}
\fancyhead[L]{\small\textbf{Yumeng He}}
\fancyhead[R]{\small Web: \href{https://heyumeng.com}{heyumeng.com} \quad Email: \href{mailto:heyumeng0928@gmail.com}{heyumeng0928@gmail.com}}
\renewcommand{\headrulewidth}{0.2pt}
\setlength{\headheight}{16pt}
\setlength{\headsep}{6pt}
\fancyfoot[C]{\small Page \thepage\ of \pageref{LastPage}}
\setlength{\footskip}{18pt}

% Slightly increase line spacing across the document
\setstretch{1.0}

% Add slight space between paragraphs (adjust here)
\setlength{\parskip}{0.3em}

\begin{document}

\vspace*{-18pt}
\begingroup
\centering
\Large\bfseries Personal Statement \par
\endgroup
% \vspace{6pt}

% Introduction
How can today's technology help families who live apart feel truly present to
one another? And what would it take to build digital twins of people, objects,
and scenes that not only look realistic but also behave correctly under touch,
motion, and control? Inspired by these questions, my research focuses on
computer graphics and vision with an emphasis on physically based simulation. By
making high-fidelity digital twins easier to create and dependable in use, I aim
to bring presence, learning, and care closer to everyday life. With these goals,
I am applying to the PhD program in Computer Science at the USC Viterbi School
of Engineering.

\textbf{Finding My Research Interests.} 
I spent real time asking myself two simple questions that many others asked me
as well: "What do you want to do?" and "What are your interests?" I first
met these questions years ago during my undergraduate study at the University of Toronto, where majors are
chosen after the first year. I explored widely: economics, statistics, and
astronomy before realizing that mathematics and computer science best match how
I think and what I want to build. That dual training still anchors me today:
math gives me structure in optimization, numerics, geometry, and CS gives me the
systems instincts to turn ideas into tools.

When I entered graduate study, I made a choice to narrow my focus and
bring rigorous code together with aesthetics and interaction through graphics
and vision. During my master's at University of Southern
California, I pursued that curiosity deliberately across computer graphics: CSCI
420 Computer Graphics, CSCI 520 Computer Animation and Simulation, and CSCI 580
3D Graphics and Rendering. I didn't just enjoy these courses; I excelled in them
--- finishing at the very top of the class --- and, more importantly, learned
how theoretical choices such as discretizations, time integration, conditioning
surface as practical tradeoffs in real code. That experience clarified my
research identity: I want to develop robust, high-fidelity physically based
simulation methods for deformable and articulated objects, so that virtual
worlds behave as convincingly as they look.

\textbf{Research in Computer Graphics.}
My research experiences at UCLA's AIVC Lab made this direction concrete. I
worked on 3D object generation with articulated parts and saw a
clear gap between visual realism and practical use. Many systems stop at
surfaces and textures, yet artists and downstream applications need models whose
parts are controllable and whose articulation can be expressed as valid URDF
rather than arbitrary segmentation. I led the end-to-end pipeline. I began with
disciplined surveying and triage, reading closely when a paper seemed near our
idea, decomposing its pipeline, mapping each step against ours, and separating
surface resemblance from true overlap. I then moved quickly from thought
experiments to working prototypes, authoring minimal URDFs from scratch,
transferring URDFs across related objects to test articulation assumptions, and
scaling from toy scripts to a full diffusion-transformer pipeline. I
experimented with lightweight adaptation such as LoRA and structure-aware
conditioning in the spirit of ControlNet, and I defined measurable criteria for
success: no self-intersection at joints under motion, consistent joint limits
under actuation, stable contact, and meshes and textures that do not destabilize
downstream use. Because I was trained as a software engineer during a year-long
co-op, I wrote modular, documented, and tested code so others could reproduce
and extend our results. This phase taught me that visual plausibility only
matters when the asset also behaves correctly under physics.

As a lighthearted aside, our project had a plot twist.
We initially set out to pursue 3D scene generation with articulated objects, but
the literature review and hands-on replication showed that articulated-object
generation itself was not yet mature enough for our goals. Many promising models
were not fully open source, and several open releases produced artifacts such as
broken topology, unstable rigging, and unreliable UVs that would not survive
downstream use. We decided to narrow the scope deliberately: first solve robust
articulated asset generation and reliable URDF translation, then return to full
scene generation. That pivot taught me how research differs from course projects.
There is no fixed instruction set, only open ideas that must be scoped,
stress-tested, and refined from zero to one. I learned how to select a tractable
entry point, design feasibility checks and ablations that reveal failure modes
early, and use those signals to steer the next iteration.

\textbf{Research in Robotics.}
In my second year of the master's program, I joined the Robotic Embedded
Systems Laboratory at USC to help design a real-to-sim pipeline for tabletop
scenes. Starting from RGB-D with known intrinsics and extrinsics, we segmented
objects, reconstructed meshes, estimated pose and scale, and imported assets
into simulation for manipulation tasks. I still remember the lab interview:
unlike earlier interviews where I mostly recited course knowledge, this time I
could truly speak pipeline --- we sketched failure modes and tradeoffs on the
spot. Drawing on prior hands-on testing, I proposed using SAM for segmentation
and Trellis for mesh reconstruction rather than building meshes from camera
point clouds, which often produce artificial artifacts and low-quality surfaces.
I also explained why we would not adopt several competing models at that stage
and compared options such as TripoSG and Hunyuan3D, clarifying the tradeoffs for
our use case. More importantly, we discussed how simulation validates grasp
feasibility and how contact, friction, and restitution parameters influence
downstream success. That moment convinced me research is the right long-term
path. Those conversations turned simulation from a classroom topic into a
decision-making tool for building robust systems.

\textbf{STEM Outreach.} 
I am committed to doing work that is useful to others. In graduate school, that
means pursuing open projects that matter rather than producing results that live
only on paper. I open-source my code and notes, write clear how-to guides, and
maintain checklists for reproducibility and fair baselines. To date I have
released seventeen projects. Each time a stranger stars a repository, I feel
encouraged not for the number itself but for the signal that someone, somewhere,
was able to build on what I made. My aim is simple and consistent with the
vision of trustworthy digital twins: tools that other people can run, reuse, and
extend.

Mentorship is the other strand of this commitment. Each year I volunteer as a
mentor, I have applied to USC's Graduate Student Mentorship and
Women in Engineering Mentorship programs to support incoming and junior students.
I still remember the first time a mentee asked what I wished I had known when I first enrolled,
it felt like speaking to my younger self. I answer candidly and
specifically---how to choose courses and find a research home, how to approach
faculty, how to set up reproducible environments and manage workloads---so that
others can avoid avoidable detours. For me, open code and open counsel are two
ways of serving the same goal: lowering the barrier for others to participate,
and building a community that turns principled ideas into working systems.

\textbf{Future Work.} \textbf{!!!!!! Need to rewrite this part !!!!!}
My decision to pursue a PhD at USC grows directly from Professor Jernej
Barbič's course in physically based simulation, which crystallized the kind of
problems I want to work on and showed me how careful modeling and numerics make
systems dependable. USC is where my preparation, interests, and mentorship fit
best. USC is uniquely compelling because it lets me pursue physically based
simulation with Professor Jernej Barbič, while also connecting to adjacent
groups that strengthen perception and generation such as Professor Yue Wang and
Professor Yajie Zhao. With Professor Barbič, I want to advance robust and
efficient simulation that respects contact, friction, joint limits, and
stability, and that can be differentiable when needed for estimation and control.
My UCLA experience taught me that some goals cannot be achieved by neural
networks alone without introducing artifacts; assets that appear watertight yet
are physically implausible still interpenetrate or violate constraints. I am
particularly interested in collision handling that guarantees non-penetration at
practical time steps and in coupling deformable and articulated dynamics under
real-world actuation. With Professor Yue Wang and Professor Yajie Zhao, I see
opportunities to connect perception, segmentation, and three-dimensional
understanding to simulation-in-the-loop pipelines so that estimated structure
and materials are verified and refined by physics rather than left as purely
visual outputs.

I hope to contribute to the Viterbi community beyond papers: reproducible
codebases, internal tooling others can build on, reading groups that bridge
graphics, vision, and robotics, and mentorship that emphasizes research hygiene
and honest ablations. I also bring a perspective shaped by both analytical rigor
and design sensibilities: I care about how objects look, move, and feel --- from
numerical stability to the quiet correctness of a door that never clips through
its frame.

\textbf{Where I See Myself.} 
I left home in middle school to study abroad, often going long stretches without
seeing my family. In a traditional household where feelings are seldom spoken,
distance felt larger than miles. That experience shapes my motivation: to make
presence travel better, and to turn digital twins from a novelty into something
that helps people feel together.

My five-year undergraduate path included a full-year co-op that taught me to
ship reliable software. I enjoy coding, and I value research that others can run,
reuse, and extend. I want my work to be used, not only cited. After the PhD, I
hope to work in a research lab building systems that bridge simulation and
perception, releasing tools the community finds as practical as they are
principled. Furthering my education at USC would bring me a step closer to that
goal by accelerating the science and engineering needed to make trustworthy
digital twins part of everyday life.

\end{document}

